% Transcription — fonds Alexandre Grothendieck, shelfmark 128, batch 1 (pages 1-20).
% Established from the facsimile archives/batches/128/batch-01.pdf, prepared
% from a copy of 128.pdf supplied by the user (PDF page k = archivists' page
% k, no cover sheet), per the skill transcribe-grothendieck. The folder is
% thirty-one pages; this is the first of two batches (pages 21-31 are batch
% 2, transcribed by another pass; its file was read for notation).
%
% Pass: Opus 5.5 (claude-opus-5-5), 2026-09-25 — first pass, unchecked against the pages by a human.
%
% Dating: the folder is « s.d. » in src/content/catalogue.ts; its inventory
% group « Descentes » (cotes 126 to 133) is dated [avant 1970], and \dating
% says both, per the skill's group rule. Nothing on these leaves dates them.
%
% Pages skipped: none. Pages summarised: none — no typescript and nothing
% by Murre or Levelt occurs in this batch. Pages transcribed: 1-20, all his,
% blue ink on white paper; pages 1, 12 carry his section headings, numbered
% by circled 1 and 2 (« Réductions via le conducteur », « Résultats
% auxiliaires sur les relations d'équivalence »). He refers to the first
% run as « N° 1 » (pp. 13, 15, 17). Pages 11, 13 and 18 are half blank.
%
% What the batch is about. Descent along a finite, not necessarily flat,
% extension A -> B = A' through a nilpotent ideal (the « conductor »).
% (1) Pages 1-11: Prop. 1 — if I ⊂ A maps isomorphically onto an ideal of
% B, the same holds in every tensor power C^n = ⊗^{n+1}_A B; Cor. 1, the
% exact sequence 0 -> I -> ⊗^E_A B -> ⊗^E_{A0} B_0 -> 0, functorial in
% E, hence 0 -> C_I -> K(B/A) -> K(B_0/A_0) -> 0 of simplicial modules;
% Cor. 2-4: H^i(B/A, M) = H^i(B_0/A_0, M_0) for i ≠ 0 and M flat, with
% 0 -> M ⊗ I -> H^0(B/A, M) -> H^0(B_0/A_0, M_0) -> 0; Cor. 5 (B_0
% faithfully flat over A_0 gives cohomological descent for flat M).
% Prop. 2: with I nilpotent, a descent datum on a flat A'-module E' is
% effective as soon as its reduction mod I is (proof for I^2 = 0 through
% the sequence of descent complexes, then induction on the nilpotency index
% with a trivial lemma on products of ideals). Cor. 1-3: Spec B -> Spec A
% is then a morphism of effective (universal, for flat base change)
% descent for flat quasi-coherent modules; globalised to X' = Spec(A') ->
% X. Prop. 3: Δ^(n)_{B/A} = Ker(⊗^n_A B -> B) ≅ Δ^(n)_{B_0/A_0}; Cor. 1-2:
% ideals in the augmentation ideal, and equivalence relations on Spec B over
% Spec A, correspond bijectively to those for B_0/A_0, compatibly with
% order, finest/coarsest, rings of invariants and effectivity.
% (2) Pages 12-18: auxiliary results on equivalence relations. Prop. 1: A'
% ⊃ A with A'/A ≅ A/m = k; exactly two equivalence relations on Spec A'
% over Spec A, exactness of A -> A' ⇉ A' ⊗ A', effective descent for flat
% modules; proof reduced to A = k, A' = k × k or k[ε]. Prop. 2: A local
% artinian, A' finite, A'' = (A' ⊗ A')/J with J an equivalence ideal, such
% that δ(x') = p_2(x') - p_1(x') generates Δ for every x' ∉ A; then
% long A'/A ≤ 1. Proof in two cases, A' non-radicial (reduced to k × ... ×
% k, n = 2) and radicial (Ω = Δ/Δ^2, derivation d, Nakayama, reduction to
% A' = k + V with V^2 = 0; the page stops at « il faut prouver que V est de
% dim. 1 »).
% (3) Pages 19-20: « Lemme 1 » (over a struck « Proposition »), the lemma
% batch 2 opens on: I of square zero, B = H^0(A'/A), C_0 = H^0(A'_0/A_0),
% hypothesis C_0 = C_0^* + Im(B -> C_0) (N.B.: true for A noetherian
% semi-local, A' finite), conclusion Ker(H^1(A'/A, G_m) -> H^1(A'_0/A_0,
% G_m)) ≅ Ker(H^1(A'/A, G_a) -> H^1(A'_0/A_0, G_a)); proof by the exact
% sequences 0 -> I C.(A'/A) -> C.(A'/A[; G_m]) -> C.(A'_0/A_0[; G_m]) -> 0
% and the cokernels of H^0 -> H^1(I C.(A'/A)); Cor. 1 (both H^1 vanish
% together when they vanish mod I).
%
% On batch 2's doubtful « Z »: here, in the proof of Lemme 1, the complex
% whose H^1 appears is plainly I C.(A'/A) — the ideal I times the cochain
% complex — which suggests that the « Z C.(A'/A, M_n) » of page 21 reads
% « I C.(A'/A, M_n) ». Not edited there.
%
% Notation fixed once for the batch: \bigotimes^E_A B for his circled ⊗ with
% the index set above; K(B/A) and C^n for his bold capital (C/K-like, read
% K where it is clear, p. 3) on pages 1-3; C^{*}_D, H^0_D for the descent
% complexes of pp. 4-6; C^{\cdot}(A'/A) on pp. 19-20 for his doubled C with
% a dot, as batch 2 does; \mathcal{A}', \mathcal{O}_X for his round
% letters (p. 8); \mathfrak{m}, \mathfrak{m}', \mathfrak{r} for his gothic
% m and r; ½ local as written. Words he underlines are \emph.
%
% Readings to check first: the displayed calculation on p. 1 (second line
% overwritten); the struck block of p. 1 and the struck Lemme 1 at the top
% of p. 2; the footnote of p. 4 and « Alors ... effectivité » there; the
% struck passage and replacement at the foot of p. 6; the marginal notes of
% pp. 14, 15 (unread) and 17; the interlinear « ceci ... » on p. 16.

\documentclass[11pt,a4paper]{article}
% Le préambule commun à toutes les transcriptions du fonds Grothendieck.
%
% Il tient en une idée : **l'incertitude se compose, elle ne se résout pas.**
% Une transcription qui lisse un mot douteux a détruit la seule chose pour
% laquelle on la faisait — un lecteur doit pouvoir savoir, à chaque endroit,
% ce qui était sur le feuillet et ce qui a été ajouté. Les six macros qui
% suivent sont donc l'appareil critique, et non de la décoration.
%
% Les mêmes commandes sont reconnues par `scripts/render.mjs`, qui produit la
% vue de lecture du site. Une macro ajoutée ici doit l'être aussi là-bas,
% faute de quoi le rendu échouera bruyamment — ce qui est voulu : un rendu
% silencieusement faux est pire qu'un rendu absent.

\ProvidesPackage{grothendieck}[2026/08/08 Transcriptions du fonds Grothendieck]

\usepackage{amsmath,amssymb,amsthm}
\usepackage{mathrsfs}
\usepackage{stmaryrd}
% Les diagrammes commutatifs sont partout dans ce fonds : c'est la langue
% même de la géométrie algébrique de Grothendieck, et une transcription qui
% les rendrait en prose ne transcrirait rien.
\usepackage{tikz-cd}
% Français par défaut — c'est la langue des feuillets — mais l'anglais est
% chargé aussi : la traduction et le résumé sont des documents anglais, et la
% typographie française (espaces avant les deux-points) y serait une faute.
% Ces documents basculent par \selectlanguage{english} en tête de corps.
\usepackage[english,french]{babel}
\usepackage{fontspec}
\usepackage{geometry}
\geometry{a4paper,margin=25mm,marginparwidth=28mm}
\usepackage{marginnote}
\usepackage{xcolor}
% ulem, not soul. `soul`'s \st and \ul are built on a character-by-character
% scan that XeLaTeX's font machinery breaks: the document fails with an
% undefined control sequence rather than degrading. ulem does the same two jobs
% and is Unicode-engine safe. `normalem` stops it hijacking \emph, which the
% transcriptions use for Grothendieck's own emphasis.
\usepackage[normalem]{ulem}
\usepackage{hyperref}
\hypersetup{colorlinks=true,linkcolor=black,urlcolor=black,pdfborder={0 0 0}}

% --- Métadonnées du lot -----------------------------------------------------
% Reprises telles quelles par le script de rendu, qui les affiche en tête de
% la vue de lecture. Elles fixent ce que le fichier prétend couvrir : sans
% elles, une transcription flotte sans point d'attache dans un fonds de seize
% mille feuillets.

\newcommand{\folder}[1]{\gdef\grothFolder{#1}}
\newcommand{\batch}[1]{\gdef\grothBatch{#1}}
\newcommand{\leaves}[2]{\gdef\grothFirst{#1}\gdef\grothLast{#2}}
\newcommand{\foldertitle}[1]{\gdef\grothFoldertitle{#1}}
\gdef\grothFolder{?}\gdef\grothBatch{?}\gdef\grothFirst{?}\gdef\grothLast{?}\gdef\grothFoldertitle{}

% --- Le résumé --------------------------------------------------------------
%
% Il ouvre la lecture modernisée, sous le titre « Résumé », et il est composé
% en petit. Ce n'est pas un ornement : c'est le seul passage du document qui
% ne soit pas des mathématiques, et un lecteur doit pouvoir le reconnaître
% avant de l'avoir lu — puis le sauter s'il connaît déjà le sujet, ou s'y
% arrêter s'il ne le connaît pas. Le filet qui le suit marque la reprise du
% corps.

\newenvironment{resume}
  {\small\setlength{\parskip}{0.45em}}
  {\par\medskip\hrule\medskip}

% --- Mots-clés ---------------------------------------------------------------
%
% En anglais, en clôture du résumé de la lecture modernisée : le vocabulaire
% moderne sous lequel chercher ce que les feuillets construisent. Le manifeste
% du site (`npm run manifest`) les extrait pour étiqueter la cote entière —
% cette ligne est la source unique des tags, il n'y a pas de fichier à côté.
\newcommand{\keywords}[1]{\par\smallskip\noindent
  {\footnotesize\textsc{Keywords} --- \textit{#1}}\par}

% --- L'appareil critique ----------------------------------------------------

% Le numéro de feuillet, en marge. C'est l'ancre qui permet de régler un
% désaccord sur une formule en regardant *un* feuillet plutôt qu'une chemise
% de six cents. Le site s'en sert aussi pour tourner les pages du fac-similé
% quand on fait défiler la transcription.
\newcommand{\leaf}[1]{%
  \marginnote{\footnotesize\color{gray!70}\textsf{#1}}%
  \label{leaf:#1}\ignorespaces}

% Illisible : on ne devine pas. Un mot inventé qui se lit comme les autres est
% le pire résultat possible.
\newcommand{\ill}{\textcolor{red!60!black}{[\,\ldots\,]}}

% Lecture douteuse : le mot est donné, et signalé comme incertain.
\newcommand{\uncertain}[1]{\uline{#1}}

% Ajout de l'éditeur — une lettre manquante, un mot sous-entendu.
\newcommand{\add}[1]{\textcolor{blue!50!black}{[#1]}}

% Biffé par Grothendieck lui-même. Ce qu'il a rayé fait partie du document :
% une rature dit souvent où la pensée a changé de direction.
\newcommand{\struck}[1]{\sout{#1}}

% Note du transcripteur, sur l'état matériel du feuillet.
\newcommand{\note}[1]{\par\smallskip{\small\color{gray!60!black}\textit{[#1]}}\par\smallskip}

% Note marginale de Grothendieck, distincte du corps du texte.
\newcommand{\marginal}[1]{\par\smallskip\noindent\hspace{1em}%
  {\small\color{gray!60!black}#1}\par\smallskip}

% --- Titre ------------------------------------------------------------------

\AtBeginDocument{%
  \begin{center}
    {\large\bfseries Fonds Alexandre Grothendieck --- cote n\textsuperscript{o}~\grothFolder}\\[2pt]
    {\itshape\grothFoldertitle}\\[4pt]
    {\small feuillets \grothFirst--\grothLast\ (lot \grothBatch)}\\[2pt]
    {\footnotesize\color{gray!60!black}Transcription établie d'après le fac-similé numérisé
     par l'Université de Montpellier.\\
     Les lectures incertaines sont soulignées, les illisibles notées [\,\ldots\,],
     les ajouts entre crochets.}
  \end{center}
  \vspace{1em}}

\endinput


\folder{128}
\batch{1}
\pages{1}{20}
\dating{s.d. — le groupe « Descentes » (126 à 133) est daté [avant 1970]}
\watermark{Édition de démonstration}
\foldertitle{Descente non plate. Modules formels. Cf Notes Murre-Levelt : notes manuscrites (s.d.)}

\begin{document}

\section*{Réductions via le conducteur}
\note{titre de sa main, en tête de la page 1, souligné et précédé d'un
« 1 » cerclé}

\page{1}

\emph{Proposition 1.} $A$ un anneau, $I$ un idéal de $A$, $B$ une
$A$-algèbre \struck{contenant $A$}, telle que l'application
$\lambda \mapsto \lambda 1_B$ de $I$ dans $B$ soit \emph{injective}, et que
son image soit \struck{\ill{}} un idéal dans $B$ [nécessairement égal à
$IB$]. Pour tout entier $n \geqslant 0$ considérons l'anneau
$C^n = \bigotimes^{n+1}_A B$, et l'application $\lambda \mapsto \lambda 1_{C^n}$
de $I$ dans $C^n$ ; \struck{Alors} cette application est injective, et son
image est un idéal.
\note{« [nécessairement égal à $IB$] » est ajouté en interligne, à
l'encre plus foncée, de même que « $C^n = \bigotimes^{n+1}_A B$ » ; il
écrit $\otimes$ cerclé, l'exposant $n+1$ au-dessus et l'indice $A$
au-dessous, et $C$ en capitale grasse}

\struck{De façon précise, soit $A_0 = A/I$, \uncertain{$B_0 = B \otimes_A A_0$}
$= B/IB = B/I1_B$, alors l'application $I \to \uncertain{C^n}$
précédente, et l'application \ill{} $C^n_0 = \bigotimes_{A_0} B_0$, donnent
lieu à une suite exacte \ill{} $0 \leftarrow C^n_0 \leftarrow C^n =
\bigotimes_A B \leftarrow I$ \ill{}}
\note{passage barré d'un trait horizontal et de plusieurs traits obliques ;
la lettre lue $C$ y ressemble à un $K$}

En effet, \ill{} \ill{}
\[
  \begin{aligned}
  \lambda (b_0 \otimes \cdots \otimes b_n) &= (\lambda b_0) \otimes b_1 \cdots \otimes b_n \\
  &= (\lambda_0 \otimes 1) \otimes b_2 \cdots \otimes b_n
  && \text{pour } \lambda_0 \in I \text{ conv.} \\
  &= 1_B \otimes \bigl[\lambda_1 (b_2 \otimes \cdots \otimes b_n)\bigr] = \\
  &= 1_B \otimes \mu (1_B \otimes \cdots \otimes 1_B)
  && \mu \in I \text{ conv.} \\
  &= \mu (1_B \otimes \cdots \otimes 1_B)
  \end{aligned}
\]
d'où le résultat par récurrence sur $n$ (trivial pour $n = 0$).
\note{la seconde ligne est surchargée (un pâté d'encre entre $\lambda$ et
$1$, un $B$ sous le produit) et sa lecture n'est pas sûre ; « d'où le
résultat par récurrence … » est écrit à gauche des deux dernières
lignes, qui ont été ajoutées après ; au-dessus et au-dessous des deux
derniers produits, des accolades comptent les facteurs : $n$, puis $n+1$}

Pour l'injectivité, on remarque que
\[
  I \to \bigotimes_A B \to B
\]
induit une injection \uncertain{i.e.}\ par hyp., elle l'\uncertain{était}
\ill{}.

\page{2}

\struck{\emph{Lemme 1.} $A$ anneau, $I$ un idéal de $A$, $B$ et $C$ deux
$A$-algèbres, telles que $B/IB$ \ill{} $A$ et $C/IC$ \ill{} \ill{}
\ill{} \ill{} $I$ \quad $A \to B$ injectif et $IB = I$, de même
$A \to C$ injectif et $IC = I$. Alors $A \to B \otimes_A C = D$ est
injectif et}
\note{encadré et barré de traits obliques ; ce lemme~1 n'est pas celui des
pages 19--20}

Comme $I \cdot 1_{C^n} = IC^n$, et $C^n/IC^n = C^n_0$
$\bigl(= \bigotimes_{A_0} B_0$, où $A_0 = A/I$, $B_0 = B/IB\bigr)$ on
trouve :

\emph{Corollaire 1.} On a une suite exacte (pour tout ens.\ fini
d'indices $E$)
\[
  0 \to I \to \bigotimes^E_A B \to \bigotimes^E_{A_0} B_0 \to 0 .
\]
Si on a une application d'ens.\ finis \add{$\sigma : E \to F$}, on a
\ill{} \uncertain{correspondants} dans le diagramme commutatif

\begin{tikzcd}[column sep=small]
0 \arrow[r] & I \arrow[r] \arrow[d, no head, "\mathrm{id}"'] & \bigotimes^E_A B \arrow[r] \arrow[d, "\sigma_{B/A}"] & \bigotimes^E_{A_0} B_0 \arrow[r] \arrow[d, "\sigma_{B_0/A_0}"] & 0 \\
0 \arrow[r] & I \arrow[r] & \bigotimes^F_A B \arrow[r] & \bigotimes^F_{A_0} B_0 \arrow[r] & 0
\end{tikzcd}

\note{« $E \to F$ » est en interligne, au-dessus de « finis » ; la
verticale de gauche est un double trait marqué « id »}

où les flèches verticales sont les homs.\ \uncertain{simpliciaux} \uncertain{déduits}
de $\sigma$, d'où une suite exacte de \add{$A$-modules}
\struck{\ill{}} simpliciaux
\[
  0 \to C_I \to C(B/A) \to C(B_0/A_0) \to 0
\]
où $C_I$ est le $A$-module simplicial \emph{\uncertain{trivial}} défini par
$I$ (\uncertain{foncteur} \struck{\ill{}} $E \mapsto I$).
\note{« $A$-modules » en interligne au-dessus d'un mot barré ; après
$C(B/A)$ un mot noirci}
\note{la lettre lue $C$ dans $C(B/A)$, comme dans $C^n$ à la page 1, est
une capitale grasse qui tient du $C$ et du $K$ ; à la page 3 il écrit
nettement $K(B/A)$}

\page{3}

\emph{Corollaire 2.} L'hom.\ $K(B/A) \to K(B_0/A_0)$ définit un
isomorphisme
\[
  H^i(B/A) \xrightarrow{\ \sim\ } H^i(B_0/A_0) \qquad \text{pour } i \neq 0
\]
et une suite exacte
\[
  0 \to I \to H^0(B/A) \to H^0(B_0/A_0) \to 0
\]
\marginal{d'où
$\mathrm{Coker}\bigl(A \to H^0(B/A)\bigr) \simeq
\mathrm{Coker}\bigl(A_0 \to H^0(B_0/A_0)\bigr)$}\note{écrit en travers,
le long de la marge gauche, et appelé par un signe d'insertion placé après
la suite exacte}

\emph{Cor.~3} a) Pour que $H^i(B/A) = 0$ il f.\ et s.\ que
$H^i(B_0/A_0) = 0$ ($i \neq 0$ donné)

b) $A \to H^0(B/A)$ inj (surj) ssi $A_0 \to H^0(B_0/A_0)$ inj (surj).
\note{« Cor.~3 », souligné, et les lettres a), b) sont ajoutés dans la
marge, devant deux lignes serrées sous la suite exacte}

Plus généralement, \uncertain{remplaçons} \ill{} \ill{} par un
$A$-module $M$ plat, on trouve :

\emph{Corollaire 4} \struck{\ill{}} Soit $M$ un $A$-module \emph{plat}.
On a une suite exacte de \add{$A$-modules simpliciaux}
\struck{complexes}
\[
  0 \to C_{M \otimes_A I} \to K(B/A, M) \to K(B_0/A_0, M_0) \to 0
\]
d'où des isomorphismes
\[
  H^i(B/A, M) \simeq H^i(B_0/A_0, M_0) \qquad i \neq 0
\]
et une suite exacte
\[
  0 \to M \otimes_A I \to H^0(B/A, M) \to H^0(B_0/A_0, M_0) \to 0 .
\]
\note{le « 4 » est surchargé sur un chiffre (3 ?) et suivi d'un mot barré ;
« plat » est souligné deux fois}

\page{4}

\emph{Corollaire 5.} Supposons $B_0$ fid.\ plat sur $A_0$. Alors pour
tout $A$-module plat \add{$M$}, on a
\[
  \begin{cases}
    H^i(B/A, M) = 0 & i \neq 0 \\
    H^0(B/A, M) \xleftarrow{\ \sim\ } M
  \end{cases}
\]

\emph{NB} Ce résultat est \emph{équivalent} au résultat correspondant dans
\uncertain{le cas} où $M = A$, i.e.\ à
\[
  \begin{cases}
    H^i(B/A) = 0 & i \neq 0 \\
    H^0(B, A) \xleftarrow{\ \sim\ } A
  \end{cases}
\]
\note{il écrit bien $H^0(B, A)$ dans la seconde accolade}

\emph{Proposition 2.} $A$, $I$, $B = A'$ comme \uncertain{ci}-dessus
\add{(tels que$^{(*)}$ $H^0(B/A) \xleftarrow{\ \sim\ } A$,
$H^1(B/A) \simeq 0$ (cf \emph{cor.~3}))}, $E'$ un $A'$-module
\emph{plat}, muni d'une donnée de descente relative à $A'/A$
[d'où $E''$, $E'''$ …\ formant un complexe
$C^{*}_D(A'/A ; E')$]. Supposons que la donnée de descente sur $E'_0$
relative à $A'_0/A_0$ déduite par changement de base $A_0 \to A'_0$ soit
effective [pour la catégorie fibrée des Modules plats], i.e.\
$\exists\, E_0$ module plat sur $A_0$, et un isom.\ compatible aux
descentes $E_0 \otimes_{A_0} A'_0 \simeq E'_0$. Alors \ill{} \ill{}
\uncertain{effectivité}
\marginal{et $I$ \emph{nilpotent}}\note{dans la marge gauche, devant
« Proposition 2 », « et $I$ nilpotent » (le dernier mot souligné deux
fois), appelé par un signe d'insertion ; au-dessus, un mot encadré
entièrement noirci de hachures. La parenthèse « tels que …\ » est
en interligne et reliée par une accolade à « dessus » ; dans
$C^{*}_D$ l'indice $D$ est ajouté}

\begin{tikzcd}[column sep=small, row sep=small]
 & E' \arrow[r, bend left=8] \arrow[r, bend right=8] \arrow[dd, no head] & E'' \arrow[r, bend left=12] \arrow[r] \arrow[r, bend right=12] & E''' \\
E_0 \arrow[r, dashed] \arrow[dd, dashed, no head] & E'_0 \arrow[r, bend left=8] \arrow[r, bend right=8] & E''_0 \arrow[r, bend left=12] \arrow[r] \arrow[r, bend right=12] & E'''_0 \\
A \arrow[r] & A' \arrow[r, bend left=8] \arrow[r, bend right=8] & A'' \arrow[r, bend left=12] \arrow[r] \arrow[r, bend right=12] & A''' \\
A_0 \arrow[r] & A'_0 \arrow[r, bend left=8] \arrow[r, bend right=8] & A''_0 \arrow[r, bend left=12] \arrow[r] \arrow[r, bend right=12] & A'''_0
\end{tikzcd}

\note{le dessin est en perspective : deux étages, $E$ au-dessus de $A$,
chacun avec sa réduction $({}_0)$ décalée en avant ; de petits traits
obliques relient chaque terme à sa réduction et des traits verticaux
chaque $E$ à son $A$ ; seuls ceux de $E'$ et de $E_0$ sont redessinés ici.
Le trait $E_0 \to E'_0$ et la verticale sous $E_0$ sont en pointillé ; les
doubles et triples traits horizontaux sont rendus par des arcs}

\marginal{$(*)$ Ces conditions \struck{\ill{}} signifient que
$\mathrm{Spec}\, B \to \mathrm{Spec}\, A$ est un morphisme de descente
\uncertain{pour} la cat.\ fib.\ des modules plats, et équivalent : la même
cond.\ \uncertain{pour} $B_0/A_0$ par \emph{cor.~3}}\note{note de bas de
page de sa main, sous un trait, appelée par l'astérisque cerclé de
l'interligne}

\page{5}

également pour la donnée de descente sur $E'$ rel\supplied{a}t\supplied{ive à}
$A'/A$, i.e.\ si $E \subset E'$ est le sous-$A$-module des
\emph{invariants} de $E'$, $E = H^0_D(A'/A ; E')$, alors $E$ est
\emph{plat} sur $A$ et $E \otimes_A A' \to E'$ est un isomorphisme.

\emph{Dém.} \marginal{Supposons d'abord $I^2 = 0$}\note{en interligne,
entouré, « Supposons » souligné} On a d'abord \uncertain{grâce à} une suite
\struck{longue de complexes} exacte de $A$-modules simpliciaux (sans
hypothèse d'effectivité pour $E'_0$ rel.\ à $A'_0/A'$)
\[
  0 \to I C^{*}_D(A'/A, E') \to C^{*}_D(A'/A, E') \to C^{*}_D(A'_0/A_0, E'_0) \to 0
\]
\note{il écrit bien $A'_0/A'$ dans la parenthèse}
\struck{$C^{*}_D(A'$} et \ill{} \ill{} pour tout $n$, \uncertain{en}
\uncertain{vertu} de la \emph{platitude} de $E'/A'$ \add{donc} de
$E^{(n)}/A^{(n)}$, et de $I^2 = 0$ :
\[
  (*) \qquad I C^n_D(A'/A, E') \simeq E^{(n)}_0 \otimes_{A^{(n)}_0}
  \underset{\substack{\| \\ I}}{I A^{(n)}} \sim
\]
$\bigl[$où $A^{(n)} = \bigotimes^{n+1}_A A'$, $E^{(n)} = E \otimes_A A^{(n)}$
\uncertain{etc}$\bigr]$. Grâce à l'effectivité pour $E'_0$ rel.\
$A'_0/A_0$, on \uncertain{peut} écrire $(*)$ sous la forme
\[
  \bigl(E_0 \otimes_{A_0} A^{(n)}_0\bigr) \otimes_{A^{(n)}_0} I A^{(n)}
  \simeq E_0 \otimes_{A_0} \bigl(\underset{\substack{\| \\ I}}{I A^{(n)}}\bigr)
  \simeq E_0 \otimes_{A_0} I
\]
d'où la \struck{\ill{}} suite exacte de $A$-modules simpliciaux :
\[
  0 \to C^{*}_{E_0 \otimes_{A_0} I} \to C^{*}_D(A'/A, E') \to C^{*}_D(A'_0/A_0, E'_0) \to 0
\]
\note{sous $E^{(n)}_0$ et sous les deux derniers $E_0 \otimes_{A_0}$ des
accolades ; dans $E^{(n)}_0$ un astérisque est surchargé ; le mot lu
« donc » est en interligne}

\page{6}

On en conclut la suite exacte
\[
  0 \to E_0 \otimes_{A_0} I \to
  \underset{\substack{\| \\ E}}{H^0_D(A'/A, E')} \to
  \underset{\substack{\| \\ E_0}}{H^0_D(A'_0/A_0, E'_0)} \to 0
\]
compte tenu que
\[
  H^i\bigl(C_{E_0 \otimes_{A_0} I}\bigr) =
  \begin{cases}
    E_0 \otimes_{A_0} I & \text{pour } i = 0 \\
    0 & \text{pour } i = 1 \\
  \end{cases}
  \quad (\text{et } \uncertain{\text{même}}\ i \neq 0)
\]
\note{les indices $D$ sous les deux $H^0$ sont ajoutés à l'encre plus
foncée}

Cela prouve d'abord que $E \to E_0$ est \emph{surjectif},
\struck{donc $\mathrm{Im}(E \to E')$ engendre le $A'$-module $E'$,} \add{puis}
et de façon précise que

a) $\mathrm{Ker}(E \to E_0) = IE$, donc $E \to E_0$ induit un
\emph{isomorphisme} $E/IE = E \otimes_A A_0 \xrightarrow{\ \sim\ } E_0$

b) $\underbracket{E/IE} \otimes_{A_0} I \to IE$ un \emph{isomorphisme},
donc grâce à a) et $E_0$ plat sur $A_0$, $E$ est \emph{plat} sur $A$.
\note{« puis » est écrit dans un petit cadre, au-dessus d'un « donc »
barré ; « plat » est souligné deux fois}

Considérons maintenant \struck{le diagramme \ill{}} \add{l'hom.}
\[
  E \otimes_A A' \to E'
\]
c'est un hom.\ \ill{} \add{de} $A'$-modules \emph{plats} qui
\struck{\ill{} \ill{} \ill{} diagramme}
\note{sous « c'est un hom. », un petit carré barré, dont on lit
$E_0 \otimes_{A_0} A'_0 \xrightarrow{\ \sim\ } E'_0$ en bas et une flèche
verticale « surj » à gauche, relié par un trait au mot « diagramme »}

\struck{\ill{} \ill{} le diagramme \ill{} \ill{} $E \otimes_A A' \to E' \to E'_0$
est surjectif, donc par Nakayama sur $A'$ (\ill{} \ill{} $IA'$
nilpotent) \ill{} \ill{} \ill{} \ill{}. \ill{} \ill{} \ill{}}
\note{deux lignes et demie barrées d'un trait et de traits obliques,
remplacées par ce qui suit}

\ill{} \ill{} \ill{} \ill{} \ill{} \ill{} \ill{} \ill{} $E \otimes_A A'$
\ill{} \ill{} $IA'$, donc est un isomorphisme.

\page{7}

Dans le cas général
\[
  I^n = 0
\]
on procède par récurrence sur $n$, \uncertain{en} notant que c'est
trivial pour $n = 1$, déjà fait pour $n = 2$, on peut supposer $n > 2$ et
le \uncertain{dém.}\ fait pour $n - 1$. \struck{Il} Il suffit alors
d'utiliser l'hyp.\ de récurrence et le

\emph{Lemme.} Si $I$, $J$ sont deux idéaux de $A$ tels que
$\lambda \mapsto \lambda 1_B$ est une bijection de $I$ ($J$) sur un idéal
de $B$, alors il en est de même de $IJ$. \uncertain{Cela} s'applique,
donc, aux puissances $I^{\nu}$ de $I$, avec $\nu \geqslant 1$.
\struck{Si de plus}

(Lemme trivial)

$+$

\emph{Corollaire 1.} Soient $A$, $I$, $B = A'$ comme dessus
[$A \subset B$, $I \subset A$ un idéal \emph{nilpotent} de $A$ tel que
$IB = I$]. \struck{et tel que} Supposons $B_0/A_0$ \struck{fid.\ plat}
\add{morphisme de} descente effective pour la cat.\ fibrée des modules
plats, p.ex.\ $B_0/A_0$ fid.\ plat. Alors
$\mathrm{Spec}\, B \to \mathrm{Spec}\, A$ est un morphisme de descente
\emph{effective} pour la catégorie fibrée des modules \emph{plats}.
\note{« morphisme de » et « $B_0/A_0$ fid.\ plat » sont en interligne,
le second au-dessus de « Alors », où il remplace un « fid.\ plat » barré
placé entre « Spec » et « $B$ »}

\page{8}

On peut \struck{\ill{}} noter le corollaire immédiat

\emph{Corollaire 2.} Soient \struck{$X$} $X$ un préschéma, $I$ un Idéal
quasi-cohérent \emph{nilpotent} sur $X$, $X' = \mathrm{Spec}(\mathcal{A}')$
un $X$-préschéma, affine sur $X$, tel que $I \to \mathcal{A}'$ soit un
isomorphisme de $I$ sur un \emph{Idéal} de $\mathcal{A}'$, et tel que
$\mathcal{A}'/I\mathcal{A}'$ soit fid.\ plat sur $\mathcal{O}_X/I$
[i.e.\ $\mathrm{Spec}(\mathcal{A}'/I\mathcal{A}') = X'_0$ fid.\ plat sur
$V(I) = X_0$]. Alors $X' \to X$ est un morphisme de descente
\emph{effective} pour la catégorie \add{fibrée} des Modules
\add{quasi-coh.}\ plats.
\note{« fibrée » et « quasi-coh. » en interligne ; il écrit
$\mathcal{A}'$ et $\mathcal{O}_X$ en lettres rondes, le second
souligné}

\struck{De \ill{}} \add{\uncertain{Il suit}}

\emph{Corollaire 3.} \add{Sous les conditions précédentes,}
\struck{Soit $A$ \ill{}} pour tout morphisme \add{plat} de préschémas
$Y \to X$, posant $Y' = X' \times_X Y$, la projection $Y' \to Y$ est un
morphisme de descente \uncertain{eff.}\ universelle pour la catégorie
fibrée des Modules quasi-coh.\ plats.
\note{« Sous les conditions précédentes » en interligne au-dessus d'un
début de phrase barré ; « plat » est écrit au crayon dans une bulle
reliée à « morphisme » ; en marge droite, $Y \to X$ puis $Y \to X$,
barrés ; le $Y$ de $Y'$ est surchargé}

\page{9}

Ce résultat s'applique en particulier lorsque $X = \mathrm{Spec}\, A$,
$X' = \mathrm{Spec}\, A'$, où $A$, $A'$ satisfont les conditions du
corollaire~1 (avec $A'_0/A_0$ fid.\ plat), quand $Y$ est un
$A$-préschéma quelconque, et $Y' = Y \otimes_A A'$.
\note{fin du paragraphe marquée d'un double trait}

\struck{\emph{Proposition 3}} Reprenons le diagramme commutatif du
Cor.~1 de prop.~1, en \struck{\ill{}} prenant $F$ réduit à un pt, on
trouve un diagramme commutatif

\begin{tikzcd}[column sep=small]
0 \arrow[r] & I \arrow[r] \arrow[d, Rightarrow, no head] & \bigotimes^n_A B \arrow[r, "\mathrm{pr}"] \arrow[d] & \bigotimes^n_{A_0} B_0 \arrow[r] \arrow[d] & 0 \\
0 \arrow[r] & I \arrow[r] & B \arrow[r] & B_0 \arrow[r] & 0
\end{tikzcd}

\note{il écrit l'exposant $n$ comme une barre au-dessus du $\otimes$
cerclé ; la verticale de gauche est un double trait terminé par une
pointe}

et donc \struck{le diagramme des noyaux \ill{}}

\emph{Proposition 3.} \add{[$A$, $I$, $B$ comme dans Prop.~1.]} Soient
$\Delta^{(n)}_{B/A} = \mathrm{Ker}\bigl(\bigotimes^n_A B \to B\bigr)$,
et définissons de même $\Delta^{(n)}_{B_0/A_0}$, alors l'hom.\
$\bigotimes^n_A B \to \bigotimes^n_{A_0} B_0$ induit un \emph{isom.}
\[
  \Delta^{(n)}_{B/A} \xrightarrow{\ \sim\ } \Delta^{(n)}_{B_0/A_0}
\]
\note{« $A$, $I$, $B$ comme dans Prop.~1. » est en interligne, entouré ;
sous « Soient », un « $0 \to I \to 0$ » (?) barré}

\page{10}

\emph{Corollaire 1.} Il y a correspondance biunivoque entre les idéaux
$J$ de $\bigotimes^n_A B$ contenus dans l'idéal d'augmentation, et les
idéaux $J^0$ de $\bigotimes^n_{A_0} B_0$ contenus dans l'idéal
d'augmentation, donnée par
\[
  \left\lbrace
  \begin{aligned}
    J &\mapsto p_n(J) \\
    p_n^{-1}(J^0) \cap \Delta^{(n)}_{B/A} &\leftarrow J^0
  \end{aligned}
  \right.
\]
Pour deux idéaux \add{$J$, $J^0$} qui se correspondent, l'hom.\ naturel
$J \to J^0$ est bijectif.
\note{« $J$, $J^0$ » en interligne ; la flèche de la seconde ligne est
un $\mapsto$ tourné vers la gauche}

Notons que la correspondance précédente est compatible avec les
relations d'ordre, \struck{et \ill{}} avec la formation de $+$ et
\uncertain{$\cap$}. D'autre part, les \struck{formations de ces}
\add{applications} précédentes (pour $n$ variable) commutent
\add{(cf Prop.~1, cor.~1)} aux opérations simpliciales, dans un sens
\uncertain{évident}.
On conclut de \struck{\ill{}} \add{ce qui précède}

\emph{Corollaire 2.} (a) Il y a correspondance biunivoque (induite par
celle du cor.~1 pour $n = 2$) entre relations d'équivalence dans
\note{« (cf Prop.~1, cor.~1) » est en interligne, relié par une accolade
à la fin de la phrase ; le « (a) » est ajouté après coup}

\page{11}

$\mathrm{Spec}\, B$ sur $\mathrm{Spec}\, A$, et celles de
$\mathrm{Spec}\, B_0$ sur $\mathrm{Spec}\, A_0$. (b) Cette correspondance
est compatible avec les relations d'ordre, \struck{\ill{} \ill{}
\ill{}} \add{en particulier, la relation} d'équivalence la plus (la
moins) fine dans \struck{plus fine. On en conclut :}
$\mathrm{Spec}\, B$ sur $\mathrm{Spec}\, A$, correspond à la relation
d'équiv.\ la plus (la moins) fine de $\mathrm{Spec}\, B_0$ sur
$\mathrm{Spec}\, A_0$.
\note{« en particulier, la relation » en interligne au-dessus d'un
passage barré}

[(c) Si $J$, $J^0$ sont des idéaux d'équivalence qui se correspondent,
et \struck{$\Delta$} $C$ ($C^0$) les \uncertain{sous-anneaux} de $B$
($B_0$) des invariants, alors $C^0 \simeq C/IC \simeq C/I$,
[\uncertain{de façon} plus précise, si $\delta : B \to D = B \otimes_A B/J$
et $\delta_0 : B_0 \to D_0$ sont les homs.\ \uncertain{obvious}, alors
…\ $\delta(b) \in ID$ ssi \struck{$b \in IB$} $\delta(b) = 0$, i.e.\
$\delta_0(b_0) = 0$ ssi $\delta(b) = 0$ i.e.\ $b_0 \in C_0$ ssi
$b \in C$], donc \uncertain{on} \ill{} \ill{} $C_0$ \uncertain{aussi}.
\uncertain{Pour} que $J$ soit effectif, il f.\ et s.\ que $J^0$ le
soit …]
\marginal{\uncertain{expliciter} \struck{\ill{}} \ill{}}\note{dans la
marge gauche, en face de $\delta_0$, deux ou trois mots dont une ligne
barrée}
\note{les lettres (b) et (c) sont cerclées ; dans $D = B \otimes_A B/J$,
le premier $B$ est surchargé. La dernière ligne de crochet est repassée à
l'encre foncée ; le reste de la page est blanc}

\section*{Résultats auxiliaires sur les relations d'équivalence}
\note{titre de sa main, en tête de la page 12, précédé d'un « 2 »
cerclé}

\page{12}

\emph{Proposition 1.} Soient $A$ un anneau, $\mathfrak{m}$ un idéal
maximal de $A$, $A'$ une $A$-algèbre, contenant $A$, tel que $A'/A$ soit
un $A$-module isomorphe à $A/\mathfrak{m} = k$. Sous ces conditions

(i) Il \uncertain{y a} dans $\mathrm{Spec}\, A'$ sur $\mathrm{Spec}\, A$
exactement deux relations d'équivalence, la plus fine
[correspondant à l'égalité dans $\mathrm{Spec}\, A'$ sur
$\mathrm{Spec}\, A$] et la moins fine [correspondant au produit fibré]

(ii) La suite
\[
  0 \to A \to A' \rightrightarrows A' \otimes_A A'
\]
est exacte, \struck{plus précis} i.e.\
$H^0(A'/A) \xleftarrow{\ \sim\ } A$, et $H^i(A'/A) = 0$ pour $i \neq 0$

(iii) $\mathrm{Spec}\, A' \to \mathrm{Spec}\, A$ est un morphisme de
descente effective pour la catégorie fibrée des Modules quasi-coh.\
plats, et reste tel\supplied{le} par toute extension de la base
$X \to \mathrm{Spec}(A)$.

\page{13}

\emph{Démonstration.} $A$, $A'$ et l'idéal $\mathfrak{m}$ satisfont les
conditions de prop.~1, de plus $A_0 = A/\mathfrak{m} = k$ est un corps,
et $A'/\mathfrak{m}A' = A'/\mathfrak{m} \simeq A'_0$ est un vectoriel de
dim.~2 sur $k$, il est en particulier fid.\ plat sur $k$. Alors (ii) et
(iii) résultent du N\textsuperscript{o}~1, Prop.~1, cor.~5, et prop.~2
\add{(cor.~1 et cor.~3)}. Pour (i), utilisant N\textsuperscript{o}~1,
Prop.~3, cor.~2, on est ramené au cas où $A = k$, donc $A'$ est une
algèbre de rang~2 sur $k$. \add{\uncertain{On peut supposer} $k$ alg.\
clos.} On doit distinguer deux cas
\note{« (cor.~1 et cor.~3) » et « On peut supposer $k$ alg.\ clos. » sont
en interligne ; « $A_0$ » porte en exposant un signe douteux ; dans
$A'/\mathfrak{m}A'$ le $\mathfrak{m}$ est surchargé}

(i) \add{$A'$ réduit, donc :} $A' \simeq k \times k$, alors la question
est purement ensembliste ; il y a exactement 2 relations d'équivalence
dans un ens.\ à 2 éléments.

(ii) $A'$ \struck{$\simeq k$} non réduit, alors $A' \simeq k[\varepsilon]$,
avec $\varepsilon^2 = 0$. (\ill{} \ill{}).
\note{« $A'$ réduit, donc : » est en interligne, entouré ; la page
s'arrête là, sans la vérification du cas (ii)}

\page{14}

\emph{Proposition 2.} Soient $A$ un anneau local artinien, $A'$ une
$A$-algèbre finie, $A'' = (A' \otimes_A A')/J$, où $J$ est un idéal
d'équivalence, \add{$\Delta = \mathrm{Ker}(A'' \to A')$.}
\struck{On suppose que}
On suppose que \struck{pour tout $x \in A'$}
\[
  A \to A' \overset{p_1}{\underset{p_2}{\rightrightarrows}} A''
\]
est une suite \emph{exacte}, et que pour tout $x' \in A' - A$, l'élément
\[
  \delta(x') = p_2(x') - p_1(x')
\]
est un \emph{générateur} de l'idéal $\Delta$. [Sous ces conditions,
$\mathrm{long}\, A'/A \leqslant 1$, et par suite (cf prop.~1, (i)) on a
$A'' \simeq A' \otimes_A A'$.
\marginal{Supposons de plus les extensions résiduelles de $A'/A$
\emph{triviales}}\note{dans la marge gauche, en oblique, appelé par un
signe d'insertion placé devant « Sous ces conditions »}
\marginal{Sans l'hypothèse $\uparrow$, il est \uncertain{vrai} que
$\mathfrak{m}A' = \mathfrak{m}$, donc on est ramené \uncertain{au cas}
$A = k$}\note{au-dessous, dans la marge, à l'encre bleue, entouré d'un
trait et relié à la ligne « Soit $\mathfrak{m}' = \mathfrak{r}(A')$ » ;
la flèche montre la note précédente ; lecture douteuse}

Soit $\mathfrak{m}' = \mathfrak{r}(A')$, distinguons deux cas

1\textsuperscript{o}) $A'$ \uncertain{n'est pas} \emph{radiciel sur $A$},
\struck{\ill{}. Alors} i.e.\
$\mathrm{Ker}(A' \otimes_A A' \to A')$ \struck{$\Delta$ \ill{} \ill{}}
n'est pas nilpotent. \emph{Comme $J$ est nilpotent}, il s'ensuit que
$\mathrm{Ker}(A'' \to A') = \Delta$ n'est pas nilpotent. Mais si
$x' \in \mathfrak{r}(A') = \mathfrak{m}'$, \struck{\ill{}} alors
$\delta(x')$ est nilpotent, donc $\delta(x') A''$ est nilpotent,
d'où $\delta(x') A'' \neq \Delta$, d'où $x' \in A$ par hyp.,
d'où
\[
  \mathfrak{m}' \subset A
\]
et comme $\mathfrak{m}' \cap A \subset \mathfrak{m} \subset
\mathfrak{m}A' \subset \mathfrak{m}'$, on a
donc
\[
  \mathfrak{m}' = \mathfrak{m} = \mathfrak{m}A'
\]
\note{« radiciel sur $A$ » et « i.e.\ $\mathrm{Ker}(A' \otimes_A A' \to A')$ »
sont en interligne au-dessus d'une ligne barrée ; dans
« $\mathfrak{m}' \cap A \subset \ldots$ », en petit à droite, un
$\mathfrak{m}A$ est surchargé en $\mathfrak{m}A'$ (?)}

\page{15}

Donc on est dans les conditions du N\textsuperscript{o}~1, avec
$A_0 = A/\mathfrak{m} = k$ un corps. \struck{Comme} Nous savons par
N\textsuperscript{o}~1 prop.~3, cor.~2 qu'on aura

\begin{tikzcd}[column sep=small, row sep=small]
A \arrow[r] \arrow[d] & A' \arrow[r, bend left=8] \arrow[r, bend right=8] \arrow[d] & A'' \arrow[d] \\
A_0 \arrow[r] & A'_0 \arrow[r, bend left=8] \arrow[r, bend right=8] & A''_0
\end{tikzcd}

\marginal{\ill{} \ill{} \ill{}}\note{trois mots au crayon, en oblique
dans la marge gauche, en face du diagramme, séparés du texte par un
trait vertical}

avec la deuxième ligne exacte ; notons qu'on aura
$A'_0 = A'/\mathfrak{m}A' = A'/\mathfrak{m}'$, $A'_0$ est \emph{réduit}.
De plus, posant
\[
  \begin{aligned}
  \Delta_0 = \mathrm{Ker}(A''_0 \to A'_0)
  &\simeq \mathrm{Ker}\bigl(A''/\mathfrak{m}A'' \to A'/\mathfrak{m}A'\bigr) , \\
  &\simeq \mathrm{Ker}\bigl(A''/\mathfrak{m}A'' \to A''/\Delta + \mathfrak{m}A''\bigr) \\
  &\simeq (\Delta + \mathfrak{m}A'')/\mathfrak{m}A''
  \simeq \Delta/\Delta \cap \mathfrak{m}A'' ,
  \end{aligned}
\]
et notant que
$\Delta \simeq \mathrm{Ker}(A'' \to A') \to \Delta_0 \simeq
\mathrm{Ker}(A''_0 \to A'_0)$ est \emph{surjectif}, on voit que
\struck{$\Delta$} l'hypothèse faite sur $A'/A$ implique que pour tout
$x'_0 \in A'_0 - A_0$, \uncertain{alors} $\delta_0(x'_0) \in A'_0$ est un
\emph{générateur} de $\Delta_0$.
\note{il écrit bien « $\delta_0(x'_0) \in A'_0$ » ; dans la dernière
ligne de la chaîne d'isomorphismes le premier $\Delta$ de
$\Delta/\Delta \cap \mathfrak{m}A''$ est surchargé ; le dernier
$\Delta_0$ l'est aussi}

Utilisant le fait que les composantes locales de $A'_0$ \uncertain{et}
[qui sont des corps] sont des extensions \emph{triviales} \add{de
$k = A_0$}, et donc
\[
  A'_0 \simeq \underbrace{k \times \cdots \times k}_{n} , \quad
  \text{avec } n \geqslant 2 ,
\]
et prenant $x'_0 = (1, 0, \ldots, 0)$ on trouve $n = 2$, OK.
\note{« de $k = A_0$ » est en interligne}

\page{16}

2\textsuperscript{o}) $A'$ \emph{est radiciel sur} $A$, i.e.\
$\mathrm{Ker}(A' \otimes_A A' \to A')$ est nilpotent, donc $\Delta$
\emph{est nilpotent}. Alors posons
\[
  \Omega = \Delta/\Delta^2
\]
et distinguons deux cas

a) $\Omega = 0$ d'où $\Delta = \Delta^2$, \struck{puisque} et
$\Delta = 0$, puisque $\Delta$ est nilpotent, donc $A'' \simeq A'$ ; donc
$A \simeq \mathrm{Ker}(A' \rightrightarrows A'') = A'$, donc
$\mathrm{long}\, A'/A = 0$.

b) $\Omega \neq 0$ ; Notons que \uncertain{dans} tous les cas, posant
pour $x' \in A'$
\[
  d(x') = \delta(x') \bmod \Delta^2 \in \Omega
  = \varphi\bigl(d_{A'/A}(x')\bigr)
\]
\uncertain{bien} où $\varphi : \Omega^1_{A'/A} \to \Omega$ est
l'application canonique, $d$ est une $A$-dérivation de $A'$ dans
$\Omega$ :
\[
  d(x'y') = x'\,dy' + y'\,dx'
\]
d'où
\[
  d(\mathfrak{m}'^2) \subset \mathfrak{m}'\Omega
\]
et comme $\Omega \neq 0$, \struck{\ill{} $\mathfrak{m}'\Omega \neq \Omega$
par Naka}\add{\uncertain{ceci} \ill{} \ill{} $\mathfrak{m}'\Omega$ ne}
peut contenir un générateur de $\Omega$ par Nakayama,
d'où d'après \struck{$\mathfrak{m}$} l'hypothèse sur $A'/A$
\[
  (*) \qquad \mathfrak{m}'^2 \subset A \quad \text{d'où }
  \mathfrak{m}'^2 \subset \mathfrak{m}
\]
Or comme $A'$ est local à ext.\ résiduelle
\note{le groupe ajouté en interligne au-dessus du passage barré n'est
lu qu'en partie}

\page{17}

triviale, \struck{on en conclut} d'où
\[
  A' = A + \mathfrak{m}'
\]
on aura
\[
  \mathfrak{m}A' = \mathfrak{m}A + \mathfrak{m}\mathfrak{m}'
  \subset \mathfrak{m} + \mathfrak{m}'^2
\]
et compte tenu de $(*)$ on aura
\[
  \mathfrak{m}A' = \mathfrak{m} ,
\]
C'est à dire on est encore dans les conditions du N\textsuperscript{o}~1,
prop.~3. On conclut comme dessus (avec les notations du
N\textsuperscript{o}~1) que la suite
\[
  \underset{\substack{\| \\ k}}{A_0} \to A'_0 \rightrightarrows A''_0
\]
est encore exacte, et satisfait \struck{encore} la condition que la
suite ci-dessus : $x'_0 \in A'_0 - A_0$ implique que $\delta(x'_0)$
engendre \struck{$\Delta$} $\Delta_0$.
\marginal{\uncertain{vérifier} \uncertain{les} \ill{}
\uncertain{préliminaires}}\note{dans la marge gauche, en oblique,
séparé du texte par un trait vertical}

On peut donc encore supposer que $A$ est un corps $k$, et de plus $A'$ est
local à idéal maximal de carré nul, donc de la forme
\[
  A' \simeq D_k(V) \simeq k + V
\]
où $V$ est un vectoriel de dimension finie sur $k$.

\page{18}

[Il faut prouver que $V$ est de dim.~1 [il ne peut être de dim.~0, car on
aurait $A' = A$ d'où $\Delta = 0$, contrairement à l'hypothèse]].
\struck{Soit $v \in$} Notons que
\[
  A' \otimes_k A' \simeq k\, 1 \otimes 1 + V \otimes 1 + 1 \otimes V
  + V \otimes V
\]
et que l'idéal engendré par
\[
  \delta(v) = v \otimes 1 - 1 \otimes v
\]
est \struck{\ill{}} \add{engendré additivement par} les multiples de
$\delta(w)$ et les $v \otimes w$ et $w \otimes v$.
\note{« engendré additivement par » en interligne au-dessus d'un mot
barré ; le crochet ouvrant du début est surchargé. La moitié inférieure
de la page est blanche}

\section*{Le lemme 1}
\note{titre de l'éditeur ; c'est le « lemme~1 » auquel renvoient les
corollaires~2 et~3 des pages 21--22}

\page{19}

\emph{\struck{Proposition} \add{Lemme 1}} Soit $A$ un anneau, $A'$ une
$A$-algèbre, $I$ un idéal \struck{nilpotent} de carré nul dans $A$,
$A_0 = A/I$, $A'_0 = A' \otimes_A A_0$. On suppose que,
\struck{$H^1(A'_0/A_0) = e$, $A_0 \xrightarrow{\ \sim\ } H^0(A'_0/A_0)$
\ill{}}
\struck{$H^1(A'_0/A_0, G_a) = 0$, $H^1(A'_0/A_0, G_m)$} posant
\[
  \begin{aligned}
  B &= H^0(A'/A) = \mathrm{Ker}(A' \rightrightarrows A'') \\
  C_0 &= H^0(A'_0/A_0) = \mathrm{Ker}(A'_0 \rightrightarrows A''_0)
  \end{aligned}
\]
on ait
\[
  C_0 = C_0^{*} + \mathrm{Im}(B \to C_0)
\]
[N.B. Cette condition est vérifiée si $A$ est noeth.\ ½ local et $A'$
fini sur $A$, car alors $\mathrm{Spec}\, C_0 \to \mathrm{Spec}\, B_0$ est
radiciel surjectif, donc $B_0$ et $C_0$ sont semi locaux et l'hom.\
$B_0 \to C_0$ induit une \emph{bijection}
$\mathrm{Max}\, C_0 \to \mathrm{Max}\, B_0$, d'où facilement l'assertion
faite]. Sous ces conditions, \uncertain{il y a} un isomorphisme canonique
\[
  (*) \quad
  \mathrm{Ker}\bigl(H^1(A'/A, G_m) \to H^1(A'_0/A_0, G_m)\bigr)
  \simeq
  \mathrm{Ker}\bigl(H^1(A'/A, G_a) \to H^1(A'_0/A_0, G_a)\bigr)
\]
\note{« Lemme 1 » est écrit au-dessus de « Proposition » barré. Après
« On suppose que, » la première ligne de conditions est barrée d'un
trait ; la seconde, encadrée, est barrée par une ligne ondulée. Le
$\simeq$ de $(*)$ est surchargé}

\emph{Démonstration.} On a les suites exactes de complexes
\[
  0 \to I C^{\cdot}(A'/A) \to C^{\cdot}(A'/A) \to C^{\cdot}(A'_0/A_0) \to 0
\]
et
\[
  0 \to I C^{\cdot}(A'/A) \to C^{\cdot}(A'/A ; G_m) \to C^{\cdot}(A'_0/A_0 ; G_m) \to 0
\]
où la deuxième est définie grâce au fait que $I^2 = 0$. Par la suite
exacte de cohomologie, on trouve donc que les deux membres de $(*)$
\note{il écrit ici un $C$ double surmonté d'un point, qu'on rend par
$C^{\cdot}$ comme dans le lot 2 ; dans la seconde suite le premier terme
est sans point. Le facteur $I$ devant $C^{\cdot}(A'/A)$ est net sur cette
page}

\page{20}

sont isomorphes respectivement à :
\[
  \mathrm{Coker}\bigl(H^0(A'_0/A_0, G_m) \to H^1(I C^{\cdot}(A'/A))\bigr)
\]
et :
\[
  \mathrm{Coker}\bigl(H^0(A'_0/A_0, G_a) \to H^1(I C^{\cdot}(A'/A))\bigr)
\]
Or on a un diagramme commutatif d'applications

\begin{tikzcd}
B^{*} \arrow[r] \arrow[d, "i"] & C_0^{*} \arrow[r] \arrow[d, "j"] & H^1(I C^{\cdot}(A'/A)) \arrow[d, "\mathrm{id}"] \\
B \arrow[r] & C_0 \arrow[r] & H^1(I C^{\cdot}(A'/A))
\end{tikzcd}

où $i$, $j$ sont les inclusions canoniques (\uncertain{celles-ci} ne sont
pas multiplicatives, mais cela est sans importance). Le résultat résulte
de cela, et de l'hypothèse
\[
  C_0 = C_0^{*} + \mathrm{Im}(B \to C_0) .
\]

\emph{Corollaire 1.} Supposons que
\[
  H^1(A'_0/A_0, G_a) = 0 \qquad H^1(A'_0/A_0, G_m) = 0
\]
\struck{Pour qu'il existe} \add{Alors il existe} un \struck{bijection}
isomorphisme can.
\[
  H^1(A'/A, G_m) \simeq H^1(A'/A, G_a)
\]
donc les relations $H^1(A'/A, G_m) = 0$ et $H^1(A'/A, G_a) = 0$ sont
équivalentes.
\note{dans $H^1(A'_0/A_0, G_a)$ l'exposant $1$ est surchargé}
\end{document}
